\chapter{Lezione 12 - 31/10}

\section{Teorema di Binet}
\Teorema{Teorema di Binet}{
    Siano \( A, B \in M_n(K) \).  
    Allora il determinante del prodotto è uguale al prodotto dei determinanti, cioè:
    \[
        \det(AB) = \det(A) \cdot \det(B).
    \]
}

\section{Secondo Teorema di Laplace}
\Teorema{Secondo Teorema di Laplace}{
    Sia \( A \in M_n(K) \) una matrice quadrata di ordine \( n \).
    Allora valgono le seguenti proprietà:

    \begin{itemize}
        \item[\textbf{(i)}] Per ogni coppia di righe distinte \( h, \bar{h} \in \{1, \dots, n\} \), si ha:
        \[
            a_{h1}A_{\bar{h}1} + a_{h2}A_{\bar{h}2} + \dots + a_{hn}A_{\bar{h}n} = 0.
        \]
        In altre parole, se si considera la matrice \( B \) ottenuta da \( A \) sostituendo la riga \(\bar{h}\) con la riga \(h\),
        allora il determinante di \(B\) è nullo:
        \[
            \det(B) = 0.
        \]

        \item[\textbf{(ii)}] Analogamente, per ogni coppia di colonne distinte \( k, \bar{k} \in \{1, \dots, n\} \), vale:
        \[
            a_{1k}A_{1\bar{k}} + a_{2k}A_{2\bar{k}} + \dots + a_{nk}A_{n\bar{k}} = 0.
        \]
        Cioè, se si considera la matrice ottenuta da \( A \) sostituendo la colonna \(\bar{k}\) con la colonna \(k\),
        il suo determinante risulta anch’esso nullo.
    \end{itemize}

    In sintesi, il teorema afferma che ogni riga (o colonna) di una matrice è ortogonale,
    rispetto ai complementi algebrici, a tutte le altre righe (o colonne) diverse da essa.

    Sia 
    \[
    A =
    \begin{pmatrix}
    1 & 2 & 3 \\
    4 & 5 & 6 \\
    7 & 8 & 9
    \end{pmatrix}
    \in M_3(\mathbb{R}).
    \]

    Scegliamo due righe distinte:
    \[
    \text{riga } h = (1, 2, 3), \quad \text{riga } \bar{h} = (4, 5, 6).
    \]

    Costruiamo ora una matrice \(B\) sostituendo la seconda riga (\(\bar{h}\)) di \(A\) con la prima riga (\(h\)):
    \[
    B =
    \begin{pmatrix}
    1 & 2 & 3 \\
    1 & 2 & 3 \\
    7 & 8 & 9
    \end{pmatrix}.
    \]

    \[
    \det(B) = 0.
    \]

}

\section{Teorema di Laplace generalizzato}
\Teorema{Teorema di Laplace generalizzato}{
Sia \( A \in M_n(K) \).  
Allora, per ogni coppia di indici \( h, \bar{h} \in \{1, \dots, n\} \) e  
per ogni coppia di indici \( k, \bar{k} \in \{1, \dots, n\} \), valgono le seguenti relazioni:

\[
    a_{h1}A_{\bar{h}1} + a_{h2}A_{\bar{h}2} + \dots + a_{hn}A_{\bar{h}n}
    = \delta_{h\bar{h}} \, |A|,
\]
\[
    a_{1k}A_{1\bar{k}} + a_{2k}A_{2\bar{k}} + \dots + a_{nk}A_{n\bar{k}}
    = \delta_{k\bar{k}} \, |A|,
\]
dove
\[
    \delta_{ij} =
    \begin{cases}
        1, & \text{se } i = j, \\
        0, & \text{se } i \neq j.
    \end{cases}
\]

In altre parole:
\begin{itemize}
    \item se \( h = \bar{h} \), si ottiene lo sviluppo di Laplace rispetto alla riga \(h\);
    \item se \( h \neq \bar{h} \), la combinazione lineare dei complementi algebrici di righe diverse è nulla;
    \item se \( k = \bar{k} \), si ottiene lo sviluppo di Laplace rispetto alla colonna \(k\);
    \item se \( k \neq \bar{k} \), la combinazione lineare dei complementi algebrici di colonne diverse è nulla.
\end{itemize}

}

\BoxBlue{Osservazione:}{}{
Se \(B\) è ridotta a gradini, allora è triangolare superiore e, di conseguenza,
\[
|B| = b_{11} \cdot b_{22} \cdot \dots \cdot b_{nn}.
\]

\medskip
Il determinante di una matrice triangolare superiore è dato dal prodotto degli elementi della sua diagonale principale.
}

\BoxBlue{Osservazione:}{}{
Il determinante di una matrice quadrata \(A \in M_n(K)\) è diverso da zero se e solo se il rango di \(A\) è massimo, cioè
\[
|A| \neq 0 \quad \Longleftrightarrow \quad \text{rango}(A) = n.
\]
}

\BoxBlue{Esercizio}{Verifica Somma Diretta}{
    Sia \(V = \mathbb{R}[x]_{\leq 3}\) uno spazio vettoriale di dimensione 4.
    Siano dati due sottospazi:
    \begin{itemize}
        \item \(W = L(w_1, w_2) = L(1+x^2, 1-x-x^2)\)
        \item \(U = L(u_1, u_2) = L(2-x, x+x^2+x^3)\)
    \end{itemize}
    Vogliamo verificare se la somma \(U+W\) è una somma diretta.

    \Proposizione{Metodi di Verifica}{
        Per definizione, la somma \(W+U\) è diretta, e si scrive \(W \oplus U\) se e solo se l'intersezione è banale:
        \[ W \cap U = \{\bar{0}\} \]
        
        Poiché \(\dim(W) = 2\) e \(\dim(U) = 2\), possiamo usare la Formula di Grassmann e il Teorema delle Basi. Le seguenti affermazioni sono equivalenti:
        \begin{itemize}
            \item \(W+U\) è somma diretta.
            \item \(\dim(W \cap U) = 0\).
            \item \(\dim(W+U) = \dim(W) + \dim(U) = 2 + 2 = 4\).
            \item L'insieme \(\mathcal{B}_{W+U} = \mathcal{B}_W \cup \mathcal{B}_U = \{w_1, w_2, u_1, u_2\}\) è una base di \(W+U\), e quindi di \(V\), cioè è linearmente indipendente.
        \end{itemize}
    }

    \Proposizione{Metodo 1: Test di Indipendenza (Determinante)}{
        Usiamo l'isomorfismo delle coordinate associato alla base canonica \(\mathcal{B} = (1, x, x^2, x^3)\):
        \[ \phi_{\mathcal{B}}: \mathbb{R}[x]_{\leq 3} \to \mathbb{R}^4 \]
        
        L'insieme \(S = \{w_1, w_2, u_1, u_2\}\) è L.I. in \(V\) se e solo se \(\phi_{\mathcal{B}}(S)\) è L.I. in \(\mathbb{R}^4\).
        Calcoliamo i vettori delle coordinate:
        \begin{itemize}
            \item \( \phi_{\mathcal{B}}(w_1) = \phi_{\mathcal{B}}(1+x^2) = (1, 0, 1, 0) \)
            \item \( \phi_{\mathcal{B}}(w_2) = \phi_{\mathcal{B}}(1-x-x^2) = (1, -1, -1, 0) \)
            \item \( \phi_{\mathcal{B}}(u_1) = \phi_{\mathcal{B}}(2-x) = (2, -1, 0, 0) \)
            \item \( \phi_{\mathcal{B}}(u_2) = \phi_{\mathcal{B}}(x+x^2+x^3) = (0, 1, 1, 1) \)
        \end{itemize}
        
        Per verificare se questi 4 vettori in \(\mathbb{R}^4\) sono L.I., controlliamo se il determinante della matrice \(A\) (che li ha come righe o colonne) è non nullo.
        \[
        A = \begin{pmatrix}
        1 & 1 & 2 & 0 \\
        0 & -1 & -1 & 1 \\
        1 & -1 & 0 & 1 \\
        0 & 0 & 0 & 1
        \end{pmatrix}
        \]
                
        
    }
}

\newpage

\BoxBlue{Esercizio}{Verifica Somma Diretta}{

    \Proposizione{Metodo 1:}{
        Sviluppiamo lungo la 4° riga (perché ha molti zeri):
        \[
        |A| = 0 - 0 + 0 + 1 \cdot (-1)^{4+4} \cdot \left| \begin{matrix}
        1 & 1 & 2 \\
        0 & -1 & -1 \\
        1 & -1 & 0
        \end{matrix} \right|
        \]
        
        Ora possiamo usare Sarrus per il minore \(M = \left| \begin{smallmatrix} 1 & 1 & 2 \\ 0 & -1 & -1 \\ 1 & -1 & 0 \end{smallmatrix} \right|\):
        \[
        |M| = (1 \cdot -1 \cdot 0) + (1 \cdot -1 \cdot 1) + (2 \cdot 0 \cdot -1) - [ (2 \cdot -1 \cdot 1) + (1 \cdot -1 \cdot -1) + (1 \cdot 0 \cdot 0) ]
        \]
        \[
        |M| = (0 - 1 + 0) - (-2 + 1 + 0) = -1 - (-1) = 0
        \]
        
        Dato che \(|M| = 0\), il determinante totale è \(|A| = 1 \cdot |M| = 0\).
        
        \textbf{Conclusione (Metodo 1):} Poiché \(\det(A) = 0\), i 4 vettori sono linearmente dipendenti.
        Quindi \(\dim(W+U) < 4\), e la somma \textbf{non è diretta}.
    }
    \Proposizione{Metodo 2: Calcolo delle Dimensioni (Grassmann)}{
        Come hai notato, possiamo calcolare la dimensione di \(W+U\) riducendo a gradini la matrice \(A\) (con i vettori in riga, per semplicità):
        \[
        A = \begin{pmatrix}
        1 & 0 & 1 & 0 \\
        1 & -1 & -1 & 0 \\
        2 & -1 & 0 & 0 \\
        0 & 1 & 1 & 1
        \end{pmatrix}
        \xrightarrow[R_3 \to R_3 - 2R_1]{R_2 \to R_2 - R_1}
        \begin{pmatrix}
        1 & 0 & 1 & 0 \\
        0 & -1 & -2 & 0 \\
        0 & -1 & -2 & 0 \\
        0 & 1 & 1 & 1
        \end{pmatrix}
        \xrightarrow[R_4 \to R_4 + R_2]{R_3 \to R_3 - R_2}
        \begin{pmatrix}
        1 & 0 & 1 & 0 \\
        0 & -1 & -2 & 0 \\
        0 & 0 & 0 & 0 \\
        0 & 0 & -1 & 1
        \end{pmatrix}
        \]
        La matrice ridotta ha 3 righe non nulle (3 pivot).
        Come hai correttamente osservato, \(\operatorname{rango}(A) = 3\), quindi \(\dim(W+U) = 3\).
        (La riga nulla ci conferma che i 4 vettori erano L.D. Infatti \(w_1+w_2 = (1+x^2) + (1-x-x^2) = 2-x = u_1\)).
        
        Usiamo la Formula di Grassmann:
        \[ \dim(W \cap U) = \dim(W) + \dim(U) - \dim(W+U) \]
        \[ \dim(W \cap U) = 2 + 2 - 3 = 1 \]
        
        \textbf{Conclusione (Metodo 2):} Poiché \(\dim(W \cap U) = 1\) (e non 0), la somma \textbf{non è diretta}.
    }
}

\newpage
\BoxBlue{Esercizio}{Verifica Somma Diretta}{
    \Proposizione{Metodo 3: Ricerca della Base dell'Intersezione}{
        Per trovare una base dell'intersezione, troviamo le equazioni cartesiane di $\phi_{\mathcal{B}}(W)$ e $\phi_{\mathcal{B}}(U)$ e le mettiamo a sistema.
        
        \textbf{Equazioni per $U = L((2, -1, 0, 0), (0, 1, 1, 1))$}
        Un vettore $x = (x_1, x_2, x_3, x_4)$ appartiene a $\phi_{\mathcal{B}}(U)$ se $\text{rango}(C_U) = 2$, dove:
        \[
        C_U = \begin{pmatrix} 2 & 0 & x_1 \\ -1 & 1 & x_2 \\ 0 & 1 & x_3 \\ 0 & 1 & x_4 \end{pmatrix}
        \xrightarrow{\text{riduzione}}
        \begin{pmatrix} 2 & 0 & x_1 \\ 0 & 2 & 2x_2 + x_1 \\ 0 & 0 & 2x_3 - (2x_2 + x_1) \\ 0 & 0 & 2x_4 - (2x_2 + x_1) \end{pmatrix}
        \]
        Imponendo le ultime due righe uguali a zero:
        \[ (S_U): \begin{cases} -x_1 - 2x_2 + 2x_3 = 0 \\ -x_1 - 2x_2 + 2x_4 = 0 \end{cases} \]
        
        \textbf{Equazioni per $W = L((1, 0, 1, 0), (1, -1, -1, 0))$}
        Un vettore $x = (x_1, x_2, x_3, x_4)$ appartiene a $\phi_{\mathcal{B}}(W)$ se $\text{rango}(C_W) = 2$, dove:
        \[
        C_W = \begin{pmatrix} 1 & 1 & x_1 \\ 0 & -1 & x_2 \\ 1 & -1 & x_3 \\ 0 & 0 & x_4 \end{pmatrix}
        \xrightarrow{\text{riduzione}}
        \begin{pmatrix} 1 & 1 & x_1 \\ 0 & -1 & x_2 \\ 0 & 0 & x_3 - x_1 - 2x_2 \\ 0 & 0 & x_4 \end{pmatrix}
        \]
        Imponendo le ultime due righe uguali a zero:
        \[ (S_W): \begin{cases} -x_1 - 2x_2 + x_3 = 0 \\ x_4 = 0 \end{cases} \]
        
        \textbf{Intersezione $W \cap U$}
        Mettiamo a sistema $S_U$ e $S_W$:
        \[ \begin{cases}
        -x_1 - 2x_2 + 2x_3 = 0 & (1) \\
        -x_1 - 2x_2 + 2x_4 = 0 & (2) \\
        -x_1 - 2x_2 + x_3 = 0 & (3) \\
        x_4 = 0 & (4)
        \end{cases} \]
        
        Sostituiamo (4) in (2):
        \[ -x_1 - 2x_2 + 0 = 0 \implies x_1 = -2x_2 \]
        
        Sostituiamo $x_1 = -2x_2$ in (3):
        \[ -(-2x_2) - 2x_2 + x_3 = 0 \implies 2x_2 - 2x_2 + x_3 = 0 \implies x_3 = 0 \]
        
        (Verifichiamo che i valori $x_1 = -2x_2, x_3=0, x_4=0$ soddisfino anche la (1): $(-(-2x_2) - 2x_2 + 2(0) = 0 \implies 2x_2 - 2x_2 = 0$. Verificato.)
        
        Le soluzioni del sistema dipendono da una variabile libera ($x_2$).
        Ponendo $x_2 = t$, la soluzione è $S = \{ (-2t, t, 0, 0) \mid t \in \mathbb{R} \}$.
        
        Una base per l'intersezione $\phi_{\mathcal{B}}(W \cap U)$ è:
        \[ L((-2, 1, 0, 0)) \]
        La dimensione è 1, confermando il risultato di Grassmann.
        
        Tornando ai polinomi, una base per $W \cap U$ è:
        \[ \phi_{\mathcal{B}}^{-1}(-2, 1, 0, 0) = -2 + x \]
    }
}


\Teorema{Invertibilità e Determinante}{
    Sia \(A \in M_n(K)\) una matrice quadrata su un campo \(K\).
    \(A\) è invertibile \(\iff \det(A) \neq 0\).

    \Dimostrazione{
        \textbf{(Direzione ``\(\Rightarrow\)'')} (Ipotesi: \(A\) è invertibile \(\implies\) Tesi: \(\det(A) \neq 0\))
        
        Per ipotesi, esiste la matrice inversa \(A^{-1} \in M_n(K)\) tale che:
        \[ A A^{-1} = I_n \]
        Calcoliamo il determinante di entrambi i membri:
        \[ \det(A A^{-1}) = \det(I_n) \]
        
        Per il \textbf{Teorema di Binet}, il determinante del prodotto è il prodotto dei determinanti. Inoltre, \(\det(I_n) = 1\).
        \[ \det(A) \cdot \det(A^{-1}) = 1 \]
        
        Poiché il prodotto dei due scalari \(\det(A)\) e \(\det(A^{-1})\) è 1 (che è diverso da 0), entrambi devono essere diversi da zero.
        In particolare, \(\det(A) \neq 0\), che è la tesi.
        (Come hai notato, questo dimostra anche la seconda tesi: \(\det(A^{-1}) = \frac{1}{\det(A)}\)).

        \bigskip
        
        \textbf{(Direzione ``\(\Leftarrow\)'')} (Ipotesi: \(\det(A) \neq 0\) \(\implies\) Tesi: \(A\) è invertibile)
        
        Dobbiamo costruire la matrice inversa. Sia \(A = (a_{ij})\).
        Definiamo \(A^\#\) come la \textbf{matrice dei complementi algebrici} (o cofattori):
        \[ A^\# = (A_{ij}) = \begin{pmatrix}
        A_{11} & A_{12} & \dots & A_{1n} \\
        \vdots & \vdots & \ddots & \vdots \\
        A_{n1} & A_{n2} & \dots & A_{nn}
        \end{pmatrix}
        \]
        Definiamo \((A^\#)^t\) come la sua trasposta, spesso chiamata \textbf{matrice aggiunta} di \(A\):
        \[ (A^\#)^t = \begin{pmatrix}
        A_{11} & A_{21} & \dots & A_{n1} \\
        \vdots & \vdots & \ddots & \vdots \\
        A_{1n} & A_{2n} & \dots & A_{nn}
        \end{pmatrix}
        \]
        Vogliamo dimostrare che \( A^{-1} = \frac{1}{\det(A)} (A^\#)^t \).
        Per farlo, calcoliamo il prodotto \( C = A \cdot (A^\#)^t \).  
        
        La matrice prodotto \(C\) è:
        \[
        C = A \cdot (A^\#)^t = \begin{pmatrix}
        \det(A) & 0 & \dots & 0 \\
        0 & \det(A) & \dots & 0 \\
        \vdots & \vdots & \ddots & \vdots \\
        0 & 0 & \dots & \det(A)
        \end{pmatrix} = \det(A) \cdot I_n
        \]
        Allo stesso modo si dimostra che \((A^\#)^t \cdot A = \det(A) \cdot I_n\).
        
        Poiché per ipotesi \(\det(A) \neq 0\), possiamo dividere per questo scalare:
        \[ A \cdot \left( \frac{1}{\det(A)} (A^\#)^t \right) = I_n \]
        \[ \left( \frac{1}{\det(A)} (A^\#)^t \right) \cdot A = I_n \]
        
        Abbiamo trovato una matrice (cioè \(B = \frac{1}{\det(A)} (A^\#)^t\)) che, moltiplicata per \(A\), dà l'identità. Questa è la definizione di matrice inversa \(A^{-1}\).
        Dunque \(A\) è invertibile.
    }
}


\Definizione{Orlato di un minore}{
Sia \(A \in M_{m \times n}(K)\).  
Sia \(M\) un minore di \(A\) di ordine \(h < \min\{m,n\}\).  
Un \textbf{orlato} di \(M\) è un altro minore \(M'\) di \(A\) di ordine \(h+1\) di cui \(M\) è una sottomatrice.

\Proposizione{Esempio}{
        Sia \(A\) la seguente matrice \(3 \times 4\):
        \[ A = \begin{pmatrix} 1 & 2 & 3 & 4 \\ 5 & 6 & 7 & 8 \\ 9 & 10 & 11 & 12 \end{pmatrix} \]
        Scegliamo un minore \(M\) di ordine \(h=2\), ad esempio quello ottenuto cancellando la riga 3 e le colonne 3 e 4:
        \[ M = \begin{pmatrix} 1 & 2 \\ 5 & 6 \end{pmatrix} \]
        
        Gli "orlati" di \(M\) sono i minori di ordine \(h+1=3\) che contengono \(M\). Si ottengono aggiungendo la riga 3 (l'unica riga che avevamo tolto) e una delle colonne che avevamo tolto (la colonna 3 o la 4):
        
        \begin{itemize}
            \item \textbf{Orlato 1} (aggiungendo riga 3 e colonna 3):
            \[ M'_1 = \begin{pmatrix} 1 & 2 & 3 \\ 5 & 6 & 7 \\ 9 & 10 & 11 \end{pmatrix} \]
            \item \textbf{Orlato 2} (aggiungendo riga 3 e colonna 4):
            \[ M'_2 = \begin{pmatrix} 1 & 2 & 4 \\ 5 & 6 & 8 \\ 9 & 10 & 12 \end{pmatrix} \]
        \end{itemize}
    }
}


\Teorema{di Kronecker (o degli Orlati)}{
    Sia \(A \in M_{m,n}(K)\) una matrice non nulla.
    
    Il rango di \(A\) è uguale a \(k\) (con \(1 \leq k \leq \min\{m,n\}\))
    \[ \iff \]
    Esiste un minore \(M\) di \(A\) di ordine \(k\) con \(\det(M) \neq 0\) e vale una delle due seguenti condizioni:
    
    \begin{enumerate}
        \item \(k = \min\{m,n\}\).
        (Cioè, \(k\) è già l'ordine massimo possibile, quindi non esistono orlati di \(M\)).
        
        \item \(k < \min\{m,n\}\) e \textbf{tutti} i minori che orlano \(M\) (cioè tutti i minori di ordine \(k+1\) che contengono \(M\)) hanno determinante nullo.
    \end{enumerate}
}



